Toyco produce y vende tres tipos de juguetes: trenes, camiones y coches. 
Cada juguete requiere un proceso de ensamblaje y uno final de empaquetado. 

Los límites semanales  en los tiempos disponibles para las dos operaciones son \textbf{4300 y 4600}, respectivamente, y los márgenes por tren, camión y coche son de \textbf{1, 2 y 5 }euros, respectivamente. 

Los tiempos que requiere cada tren un cada proceso (ensamble y empaquetado) son \textbf{2, 1} minutos, respectivamente. Los de cada camión son \textbf{2 y 2} y los de cada coche \textbf{1 y 4} minutos.

Además, existe un compromiso comercial con un distribuidor de juguetes de suministrarle al menos \textbf{1000} unidades entre camiones y coches.

Si $x_1$, $x_2$ y $x_3$ representan el número semanal de unidades producidas de trenes, camiones y coches, respectivamente, el modelo de programación lineal que permite establecer la producción con mayor margen es el siguiente.

\begin{equation}
\begin{split}
\mbox{max. } z =  x_1 + 2x_2 + 5x_3\\
s.a.:\\
2x_1 + 2x_2 +  x_3 \leq 4300\\
 x_1 + 2x_2 + 4x_3 \leq 4600\\
        x_2 +  x_3 \geq 1000\\
x_1,\,\,x_2,\,\,x_3\geq 0\\
\end{split}
\end{equation}

\begin{center}
\begin{tabular}{c|c|cccccc|}
 & $z$ & $x_1$ & $x_2$ & $x_3$ & $h_1$ & $h_2$ & $h_3$\\ 
Phase 2 & -5750 & $-1/4$ & $-1/2$ & 0 & 0 & $-5/4$ & 0 \\ 
\hline
$h_1$ & 3150  & $7/4$ & $3/2$ & 0 & 1 & $-1/4$ & 0\\ 
$h_3$ & 150  & $1/4$ & $-1/2$ & 0 & 0 & $1/4$ & 1\\ 
$x_3$ & 1150  & $1/4$ & $1/2$ & 1 & 0 & $1/4$ & 0\\ 
\hline
\end{tabular}
\end{center}





Se pide:

\begin{enumerate}
    \item Interpretar la información que ofrece la tabla de la solución óptima (producción, margen total y recusos y compromiso comercial).
    
    \item Calcular dentro de qué rango de valores del tiempo disponible de empaquetado la gama de productos del plan óptima de producción no cambiaría.
    
    \item Explicar cuál sería el impacto de relajar el compromiso comercial adquirido

    \item Indicar y justificar cómo afectaría esto al plan óptimo de producción si, por política de empresa, hubiera que que fabricar al menos tantos camiones como coches.
   
\end{enumerate}







